\section{Introduction}
\label{sec:introduction}
\IEEEPARstart{I}{n} contemporary smart manufacturing encompassing advanced electronics, precision machining, pharmaceutical packaging, and technical textiles, automated surface defect inspection is a critical determinant of product yield and operational profitability \cite{bergmann2019mvtec, czimmermann2020visual}. Modern vision inspection systems are transitioning beyond coarse image-level pass/fail classification toward fine-grained, pixel-level polygonal instance segmentation to support downstream robotic rejection, laser-guided repair, and root-cause process diagnosis. In this regime, inspection algorithms face a sharp economic trade-off between the \textit{Defect Escape Rate} (False Negatives), which risks shipping compromised parts and incurring costly warranty claims, and the \textit{False Alarm Rate} (False Positives), which erroneously scraps flawless goods and halts automated lines.

Despite the rapid evolution of deep convolutional architectures and real-time segmentation heads (e.g., YOLOv11-seg \cite{ultralytics2024yolo11}), supervised instance segmentation requires extensive libraries of high-resolution images with labor-intensive, pixel-accurate polygon annotations. In mature fabrication facilities, however, production yields routinely exceed $99\%$, rendering genuine defective samples exceedingly scarce and unpredictable. To overcome this data scarcity barrier, recent research has pivoted toward self-supervised anomaly synthesis paradigms---such as CutPaste \cite{li2021cutpaste}, DRAEM \cite{zavrtanik2021draem}, and Natural Synthetic Anomalies (NSA) \cite{schluter2022natural}---which synthesize defective training patterns strictly from nominal reference images.

Nevertheless, transitioning self-supervised anomaly generation from academic benchmarks to practical industrial inspection reveals three critical image-processing bottlenecks:
\begin{enumerate}
    \item \textbf{Background Lattice Domination and Boundary Artifacts in Spatial Synthesis:} Conventional anomaly synthesizers operate almost exclusively in the spatial pixel domain, creating a sharp dichotomy across material classes. On periodic textures (e.g., woven \texttt{carpet} or metallic \texttt{grid}), repetitive spatial lattices concentrate intense spectral energy that severely masks subtle synthetic defects, causing models to miss micro-cracks or overfit to weaving patterns. Conversely, on non-periodic rigid objects (e.g., glass \texttt{bottle} contours or metal \texttt{screw} threads), unconstrained patch pasting introduces sharp boundary artifacts or spills defects onto non-object backgrounds.
    \item \textbf{The Severe False-Alarm Dilemma:} Because conventional self-supervised models are trained exclusively on artificially perturbed scenes, they become hyper-sensitive to benign surface grain, minor reflections, or yarn displacements. Without an explicit mechanism to preserve confidence on nominal textures, these models suffer elevated false alarm rates that trigger costly line stoppages.
    \item \textbf{Vulnerability to Harsh Optical Degradations:} Factory cameras operate under conveyor vibration, thermal lens defocus, and low-light sensor noise. Purely spatial detectors degrade sharply under such optical corruptions---as demonstrated by common image-corruption benchmarks such as ImageNet-C \cite{hendrycks2019robustness} and its industrial adaptation MVTec-AD-C \cite{chen2023mvtecadc}---because spatial convolutions alone struggle to decouple high-frequency defect edges from low-frequency optical blur. Moreover, heavy generative or memory-bank detectors (e.g., PaDiM \cite{defard2021padim}, PatchCore \cite{roth2022towards}) output coarse heatmaps rather than polygonal masks and cannot satisfy high-throughput inline speeds.
\end{enumerate}

To simultaneously resolve these inter-domain challenges, this paper presents a unified frequency-spatial image processing framework for industrial surface defect segmentation under a \textit{Strict Defect-Free Training Regime} (zero real defective samples observed during training). On the frequency branch, we first establish coordinate-preserving invertibility and then deploy an adaptive 2D Fast Fourier Transform (2D-FFT) Gaussian notch filter to suppress periodic background lattice energy on structured textures, coupled with a 2D Discrete Wavelet Transform (2D-DWT) Haar decomposition that enhances high-frequency boundary gradients against optical defocus. On the spatial branch, we resolve the geometric breakdown of multi-octave Perlin noise on non-power-of-two industrial resolutions and enforce seamless boundary blending via foreground-clamped Poisson partial differential equations. By enforcing an explicit negative-sample preservation invariant during polygonal segmentation optimization, our framework trains the network to suppress false activations on defect-free surfaces while isolating genuine micro-defects.

The core scientific contributions of this paper are summarized as follows:
\begin{itemize}
    \item \textbf{Coordinate-Preserving Frequency-Spatial Synthesis:} We first validate the unfiltered identity roundtrip ($\mathcal{T}^{-1}(\mathcal{T}(I))$ without spectral filtering) across 2,676 MVTec AD RGB images, confirming high numerical fidelity ($\text{PSNR} = 51.31\text{ dB}$, $\text{SSIM} = 0.9983$ for 2D-FFT; exact integer identity for Haar 2D-DWT) and zero spatial coordinate shift of segmentation masks. Building on this coordinate-preserving transform backbone, we apply adaptive 2D-FFT periodic notch filtering and Haar 2D-DWT sub-band modulation to suppress periodic lattices and enhance defect boundaries.
    \item \textbf{Arbitrary-Resolution Spatial Anomaly Engineering:} We eliminate dimensional broadcast failures of multi-octave Perlin noise synthesis on non-power-of-two industrial images (e.g., $900 \times 900$ for \texttt{bottle}, $840 \times 840$ for \texttt{tile}) via canonical $512 \times 512$ scalar grid generation followed by adaptive bilinear interpolation, foreground-clamped Poisson seamless cloning, and polygonal vectorization.
    \item \textbf{Negative-Sample False-Alarm Suppression and Quantitative Benchmarking:} Under a strict defect-free training regime, we introduce an explicit negative-sample preservation invariant ($\rho_{\text{neg}} = 0.15$) paired with a dual-domain evaluation engine and MVTec-AD-C optical stress testing. Supervised YOLOv11-seg oracle baselines achieve $0.854$--$0.873$ Mask mAP@50 across six MVTec AD categories while exposing $4.50\%$--$8.78\%$ nominal $\text{FAR@0.25}$, and TensorRT FP16 model inference profiling confirms real-time feasibility at $6.7\text{ ms}$ ($\approx 149.2\text{ FPS}$) on $1024 \times 1024$ frames.
\end{itemize}

The remainder of this paper is organized as follows: Section~\ref{sec:related_work} reviews related literature across anomaly detection, self-supervised synthesis, and frequency-domain visual inspection. Section~\ref{sec:methodology} delineates the proposed inter-domain frequency-spatial synthesis architecture, negative-sample preservation, and polygonal segmentation optimization. Section~\ref{sec:experiments} presents the experimental setup, dual-domain evaluation metrics, frequency roundtrip verification, segmentation benchmarks, MVTec-AD-C optical stress testing, and runtime efficiency analysis. Finally, Section~\ref{sec:conclusion} concludes the paper and outlines future research directions.